%! Intercepts, Zeros, and Extrema — Unit 2, Lesson 2.4 — Slides (Beamer)
\documentclass[aspectratio=169, 11pt]{beamer}
\usepackage{atda-beamer}   % provides \royalheader{} and \sectionlabel[color]{}
\usepackage{pgfplots}      % atda-beamer does NOT load pgfplots
\pgfplotsset{compat=1.18}

\begin{document}

% TITLE SLIDE (CourseName is not defined in beamer, so the name is literal)
{
\setbeamercolor{background canvas}{bg=royal}
\begin{frame}[plain]
  \vspace{2.8cm}\hspace{0.55cm}%
  \begin{minipage}{0.9\textwidth}
    {\color{goldacc}\bfseries\small LESSON 2.4}\\[0.5cm]
    {\color{white}\bfseries\LARGE Intercepts, Zeros, and Extrema}\\[1.2cm]
    {\color{mistmid}\small Algebra, Trigonometry, and Data Analysis~$\cdot$~Unit 2}
  \end{minipage}
\end{frame}
}

% ── HOOK ──────────────────────────────────────────────────────────────────────
\begin{frame}[t, plain]
  \royalheader{Two Students, One Question, Both Answers in Dollars}
  \sectionlabel[goldacc]{hook --- who answered the question that was asked?}
  Last night's yearbook curve: profit is \$0 at prices of \$0 and \$60, and it peaks at \$1800 when
  the price is \$30. \textbf{The question: over which prices is the profit going up?}
  \begin{block}{Two answers, and neither student made an arithmetic mistake}
    \centering\large
    Student A: \textbf{``from 0 to 1800''} \qquad\qquad Student B: \textbf{``from 0 to 30''}
  \end{block}
  \vspace{0.15cm}
  Every number on both sides is really on that graph. \textbf{And you cannot settle it by checking
  units --- all four are dollars.}
\end{frame}

% ── THE AXIS SETTLES IT ───────────────────────────────────────────────────────
\begin{frame}[t, plain]
  \royalheader{The Axis Settles It}
  \sectionlabel{before you write a number down, say which axis it came from}
  \begin{block}{The question said \emph{over which prices} --- and a price is an input}
    \centering\large
    \textbf{Student B is right: increasing on $(0,30)$, an interval of \emph{prices}.}
  \end{block}
  \begin{itemize}
    \setlength{\itemsep}{0.3cm}
    \item A's $0$ and $1800$ are \textbf{profits} --- they say how \emph{high} the graph climbed,
          not \emph{where}.
    \item ``Increasing,'' ``decreasing,'' ``constant'' are always answers about \textbf{inputs}.
    \item ``Maximum'' and ``minimum'' are \textbf{outputs} --- and they come with an input attached.
  \end{itemize}
\end{frame}

% ── INTERCEPTS AND ZEROS ──────────────────────────────────────────────────────
\begin{frame}[t, plain]
  \royalheader{Intercepts Are Points. Zeros Are Numbers.}
  \sectionlabel{\texttt{AFDA.AF.2d} --- both are questions about a coordinate that is zero}
  \renewcommand{\arraystretch}{1.3}
  \small
  \begin{tabularx}{\textwidth}{@{}p{2.6cm} p{4.4cm} X@{}}
    \textbf{Name} & \textbf{What it is} & \textbf{How many a function can have} \\
    \hline
    $y$-intercept & the \textbf{point} $(0,b)$ where the graph crosses the vertical axis &
      at most one --- one input, one output \\
    $x$-intercept & a \textbf{point} $(k,0)$ where the graph crosses the horizontal axis &
      none, one, or many \\
    zero & the \textbf{number} $k$ in that first slot: $f(k)=0$ &
      one for each $x$-intercept \\
  \end{tabularx}
  \vspace{0.3cm}

  \textbf{Three sentences that say the same thing:} \quad $k$ is a zero of $f$ \quad$\cdot$\quad
  $f(k)=0$ \quad$\cdot$\quad the point $(k,0)$ is on the graph.
\end{frame}

% ── NOT EVERY GRAPH HAS A ZERO ────────────────────────────────────────────────
\begin{frame}[t, plain]
  \royalheader{Not Every Graph Crosses the $x$-Axis}
  \sectionlabel[goldacc]{and ``$3$'' can be an input on one graph and an output on the next}
  \begin{center}
  \begin{tabular}{@{}c@{\hspace{0.9cm}}c@{\hspace{0.9cm}}c@{}}
  \begin{tikzpicture}
  \begin{axis}[width=4.6cm, height=3.6cm, title={\tiny $y=x^2-4$},
      xmin=-3.4, xmax=3.4, ymin=-5.4, ymax=5.4,
      xtick={-2,0,2}, ytick={-4,0,4}, minor x tick num=1, minor y tick num=1,
      grid=both, grid style={linegray!45}, minor grid style={linegray!30},
      axis lines=left, tick label style={font=\tiny}]
    \addplot[royal, very thick, domain=-3:3, samples=90]{x^2-4};
  \end{axis}
  \end{tikzpicture}
  &
  \begin{tikzpicture}
  \begin{axis}[width=4.6cm, height=3.6cm, title={\tiny $y=2^x$},
      xmin=-3.4, xmax=3.4, ymin=-1.4, ymax=8.4,
      xtick={-2,0,2}, ytick={0,4,8}, minor x tick num=1, minor y tick num=1,
      grid=both, grid style={linegray!45}, minor grid style={linegray!30},
      axis lines=left, tick label style={font=\tiny}]
    \addplot[royal, very thick, domain=-3.3:3, samples=90]{2^x};
  \end{axis}
  \end{tikzpicture}
  &
  \begin{tikzpicture}
  \begin{axis}[width=4.6cm, height=3.6cm, title={\tiny $y=3-x$},
      xmin=-1.4, xmax=5.4, ymin=-2.4, ymax=4.4,
      xtick={0,2,4}, ytick={-2,0,2,4}, minor x tick num=1, minor y tick num=1,
      grid=both, grid style={linegray!45}, minor grid style={linegray!30},
      axis lines=left, tick label style={font=\tiny}]
    \addplot[royal, very thick, domain=-1.3:5.3, samples=2]{3-x};
  \end{axis}
  \end{tikzpicture}
  \end{tabular}
  \end{center}
  \vspace{-0.25cm}
  Zeros $-2$ and $2$; $y$-intercept $(0,-4)$. \quad \textbf{No zeros at all}; $y$-intercept $(0,1)$.
  \quad Zero $3$; $y$-intercept $(0,3)$ --- \textbf{same number, different axes}.
\end{frame}

% ── INCREASING / DECREASING / CONSTANT ────────────────────────────────────────
\begin{frame}[t, plain]
  \royalheader{Rising, Falling, Holding Steady --- Always Left to Right}
  \sectionlabel{\texttt{AFDA.AF.2b} --- twelve hours of temperature, midnight to noon}
  \begin{center}
  \begin{tikzpicture}
  \begin{axis}[
      width=10.4cm, height=4.2cm,
      xlabel={$t$ (hours after midnight)}, ylabel={$T$ (degrees C)},
      xmin=0, xmax=12, ymin=-6.4, ymax=6.4,
      xtick={0,2,4,6,8,10,12}, ytick={-6,-4,-2,0,2,4,6},
      minor x tick num=1, minor y tick num=1,
      grid=both, grid style={linegray!45}, minor grid style={linegray!30},
      axis lines=left, tick label style={font=\tiny}, label style={font=\tiny}]
    \addplot[royal, very thick] coordinates {(0,2) (2,2) (5,-4) (9,4) (12,1)};
  \end{axis}
  \end{tikzpicture}
  \end{center}
  \vspace{-0.3cm}
  Constant on $(0,2)$. \quad Decreasing on $(2,5)$. \quad Increasing on $(5,9)$. \quad Decreasing on
  $(9,12)$.\\
  \textbf{Every one of those is a set of \emph{times}} --- and the zeros at $t=3$ and $t=7$ say it
  was below freezing from 3 a.m.\ to 7 a.m.
\end{frame}

% ── THE THREE RULES ───────────────────────────────────────────────────────────
\begin{frame}[t, plain]
  \royalheader{Three Rules for Intervals}
  \sectionlabel{these decide every increasing/decreasing question you will be asked}
  \begin{enumerate}
    \setlength{\itemsep}{0.35cm}
    \item \textbf{Report the interval on the input axis.} Never on the output axis. This is the hook,
          and it is the single most common mistake in this lesson.
    \item \textbf{Parentheses at the ends.} At a \textbf{turning point} the graph is neither rising
          nor falling, so it belongs to neither interval.
    \item \textbf{An exponential never changes its mind.} $y=2^x$ rises on all of
          $(-\infty,\infty)$; $y=\left(\tfrac12\right)^x$ falls on all of it. One interval, every
          time.
  \end{enumerate}
\end{frame}

% ── ABSOLUTE EXTREMES ─────────────────────────────────────────────────────────
\begin{frame}[t, plain]
  \royalheader{A Value \emph{and} a Location}
  \sectionlabel[goldacc]{\texttt{AFDA.AF.2c} --- ``4 degrees'' is half an answer}
  \begin{block}{The absolute maximum and minimum are the highest and lowest outputs over the
    \emph{entire} domain}
    \centering\large
    \textbf{``$4$ degrees, at $t=9$''} --- the value from the vertical axis, the location from the
    horizontal one.
  \end{block}
  \begin{itemize}
    \setlength{\itemsep}{0.28cm}
    \item Two different axes in one sentence. That is exactly why this is where people slip.
    \item \textbf{``Absolute'' means over the whole domain} --- a high point in the middle is not
          automatically the highest.
    \item One extreme \emph{value} can happen at \emph{two} locations. Give both.
  \end{itemize}
\end{frame}

% ── THE DOMAIN DECIDES ────────────────────────────────────────────────────────
\begin{frame}[t, plain]
  \royalheader{The Domain Is Part of the Question}
  \sectionlabel{same rule $y=x^2$, two domains, different answers}
  \begin{center}
  \begin{tabular}{@{}c@{\hspace{1.2cm}}c@{}}
  \begin{tikzpicture}
  \begin{axis}[width=5.6cm, height=3.9cm, title={\tiny all real numbers},
      xmin=-3.4, xmax=3.4, ymin=-1.0, ymax=10.4,
      xtick={-2,0,2}, ytick={0,4,8}, minor x tick num=1, minor y tick num=1,
      grid=both, grid style={linegray!45}, minor grid style={linegray!30},
      axis lines=left, tick label style={font=\tiny}]
    \addplot[royal, very thick, domain=-3.15:3.15, samples=90]{x^2};
  \end{axis}
  \end{tikzpicture}
  &
  \begin{tikzpicture}
  \begin{axis}[width=5.6cm, height=3.9cm, title={\tiny domain $-1 \le x \le 3$},
      xmin=-3.4, xmax=3.4, ymin=-1.0, ymax=10.4,
      xtick={-2,0,2}, ytick={0,4,8}, minor x tick num=1, minor y tick num=1,
      grid=both, grid style={linegray!45}, minor grid style={linegray!30},
      axis lines=left, tick label style={font=\tiny}]
    \addplot[royal, very thick, domain=-1:3, samples=90]{x^2};
    \addplot[goldacc, only marks, mark=*, mark size=1.7pt] coordinates {(-1,1) (3,9)};
  \end{axis}
  \end{tikzpicture}
  \end{tabular}
  \end{center}
  \vspace{-0.3cm}
  Minimum $0$ at $x=0$, \textbf{no maximum}. \qquad Minimum $0$ at $x=0$, \textbf{maximum $9$ at
  $x=3$} --- an \emph{endpoint}.\\
  \textbf{Not one point of the curve moved.} Cutting the domain handed it an extreme it did not have.
\end{frame}

% ── GUIDED PRACTICE ───────────────────────────────────────────────────────────
\begin{frame}[t, plain]
  \royalheader{Guided Practice: The Greenhouse Cistern}
  \sectionlabel{one graph, five characteristics, every one of them interpreted}
  \begin{center}
  \begin{tikzpicture}
  \begin{axis}[
      width=9.6cm, height=4.0cm,
      xlabel={$t$ (days after Monday morning)}, ylabel={$W$ (depth, feet)},
      xmin=0, xmax=7, ymin=0, ymax=8,
      xtick={0,1,2,3,4,5,6,7}, ytick={0,2,4,6,8}, minor y tick num=1,
      grid=both, grid style={linegray!45}, minor grid style={linegray!30},
      axis lines=left, tick label style={font=\tiny}, label style={font=\tiny}]
    \addplot[royal, very thick] coordinates {(0,6) (3,0) (5,4) (6,4) (7,7)};
  \end{axis}
  \end{tikzpicture}
  \end{center}
  \vspace{-0.3cm}
  $y$-intercept $(0,6)$. \quad Zero at $t=3$ --- \textbf{the cistern ran dry on Thursday}. \quad
  Decreasing $(0,3)$, increasing $(3,5)$ and $(6,7)$, constant $(5,6)$.\\
  \textbf{Maximum $7$ ft at $t=7$; minimum $0$ ft at $t=3$.} So ``the water was highest at the start
  of the week'' is false.
\end{frame}

% ── ACTIVITY LAUNCH ───────────────────────────────────────────────────────────
\begin{frame}[t, plain]
  \royalheader{Group Activity: Reading a Graph for Everything It Says}
  \sectionlabel{groups of three --- everyone starts at Tier R}
  No equations today. Every answer is read off a picture.
  \begin{block}{The rule for the whole activity}
    Every time you write a number, say out loud whether it came from the \textbf{horizontal} axis or
    the \textbf{vertical} one. \emph{An interval of inputs and a stretch of outputs can be the same
    two numbers.}
  \end{block}
  \vspace{0.15cm}
  Tier E: two claims that read real numbers off real graphs and still answer the wrong question, one
  domain cut that changes three things at once, and a graph with no zeros at all.
\end{frame}

% ── CLOSE ─────────────────────────────────────────────────────────────────────
\begin{frame}[t, plain]
  \royalheader{Exit Ticket}
  \sectionlabel[goldacc]{notes closed --- 5 minutes}
  \begin{enumerate}
    \setlength{\itemsep}{0.3cm}
    \item The creek graph: $y$-intercept, both zeros and what a zero means, and the two intervals.
    \item Absolute maximum and minimum --- \textbf{value and location} --- and turning point or
          endpoint for each.
    \item Somebody wrote ``increasing from $-3$ to $9$.'' Those are your own item-2 numbers. Say what
          is wrong anyway.
  \end{enumerate}
  \vspace{0.25cm}
  {\color{royal}\bfseries Next: Lesson 2.5 --- End Behavior and Asymptotes.}
\end{frame}

\end{document}
